% ==============================================================
% Remarkable Pro Move briefing preamble
% --------------------------------------------------------------
% - Page: 4.00in x 7.11in (16:9 portrait, approximates Move display)
% - Body: TeX Gyre Pagella (Palatino-like Latin-script serif)
% - Headings: TeX Gyre Adventor (companion sans-serif)
% - Footnotes per page + consolidated References list at end
%
% USAGE:
%   % ==============================================================
% Remarkable Pro Move briefing preamble
% --------------------------------------------------------------
% - Page: 4.00in x 7.11in (16:9 portrait, approximates Move display)
% - Body: TeX Gyre Pagella (Palatino-like Latin-script serif)
% - Headings: TeX Gyre Adventor (companion sans-serif)
% - Footnotes per page + consolidated References list at end
%
% USAGE:
%   % ==============================================================
% Remarkable Pro Move briefing preamble
% --------------------------------------------------------------
% - Page: 4.00in x 7.11in (16:9 portrait, approximates Move display)
% - Body: TeX Gyre Pagella (Palatino-like Latin-script serif)
% - Headings: TeX Gyre Adventor (companion sans-serif)
% - Footnotes per page + consolidated References list at end
%
% USAGE:
%   % ==============================================================
% Remarkable Pro Move briefing preamble
% --------------------------------------------------------------
% - Page: 4.00in x 7.11in (16:9 portrait, approximates Move display)
% - Body: TeX Gyre Pagella (Palatino-like Latin-script serif)
% - Headings: TeX Gyre Adventor (companion sans-serif)
% - Footnotes per page + consolidated References list at end
%
% USAGE:
%   \input{remarkable-preamble.tex}
%   \newcommand{\briefingtitle}{My Briefing}
%   \newcommand{\transcriptnum}{1}
%   \title{...}
%   \author{...}
%   \date{...}
%   \begin{document}\maketitle\thispagestyle{plain}
%   ... body ...
%   \end{document}
%
% ANNOTATION MARGIN:
%   The default right margin is 0.75in, which reserves space for
%   handwritten stylus annotations on e-ink readers. This gives a
%   text column of ~2.90in — a good balance between readability and
%   annotation space. For maximum words-per-page (no annotations),
%   change to `right=0.35in`. For a wider annotation column, use
%   `right=1.00in` (~2.65in text column, ~25% more pages).
% ==============================================================

\documentclass[10pt]{article}

% ---------- Geometry (Remarkable Pro Move, ~4in wide, 16:9) ----------
\usepackage[
  paperwidth=4.00in,
  paperheight=7.11in,
  top=0.40in,
  bottom=0.45in,
  left=0.35in,
  right=0.75in,         % annotation margin for stylus notes (0.35in for no margin)
  footskip=0.20in,
  headheight=12pt,
  headsep=0.10in
]{geometry}

% ---------- Fonts ----------
\usepackage[T1]{fontenc}
\usepackage[utf8]{inputenc}
\usepackage{tgpagella}              % Pagella for body (Latin-script serif)
\renewcommand{\familydefault}{\rmdefault}
\usepackage{tgadventor}             % Adventor (Latin-script sans) companion
\renewcommand{\sfdefault}{qag}      % set sans default to Adventor
\usepackage{microtype}              % cleaner line breaks and spacing

% ---------- Layout helpers ----------
\usepackage{setspace}
\setstretch{1.14}                   % mild leading for e-ink readability
\usepackage{parskip}                % blank line between paragraphs, no indent
\setlength{\parskip}{0.55em plus 0.15em minus 0.10em}

% narrow-column line breaking: allow TeX more slack before overfulling
\setlength{\emergencystretch}{3.5em}
\hyphenpenalty=50
\tolerance=2000

% ---------- Sectioning ----------
\usepackage{titlesec}
% Smaller headings — page is only 4in wide
\titleformat{\section}{\sffamily\large\bfseries}{\thesection}{0.6em}{}
\titleformat{\subsection}{\sffamily\normalsize\bfseries}{\thesubsection}{0.5em}{}
\titleformat{\subsubsection}{\sffamily\small\bfseries\itshape}{}{0em}{}
\titlespacing*{\section}{0pt}{1.4em}{0.6em}
\titlespacing*{\subsection}{0pt}{1.1em}{0.45em}
\titlespacing*{\subsubsection}{0pt}{0.8em}{0.30em}

% ---------- Page headers/footers ----------
\usepackage{fancyhdr}
\pagestyle{fancy}
\fancyhf{}
\renewcommand{\headrulewidth}{0.3pt}
\renewcommand{\footrulewidth}{0pt}
% The document MUST define \briefingtitle and \transcriptnum with
% \newcommand BEFORE \begin{document} (see example-body.tex).
\fancyhead[L]{\scriptsize\sffamily\MakeUppercase{\briefingtitle}}
\fancyhead[R]{\scriptsize\sffamily \transcriptnum}
\fancyfoot[C]{\small\thepage}

% Plain style for first page
\fancypagestyle{plain}{%
  \fancyhf{}%
  \renewcommand{\headrulewidth}{0pt}%
  \fancyfoot[C]{\small\thepage}%
}

% ---------- Lists ----------
\usepackage{enumitem}
\setlist{nosep, leftmargin=1.1em, itemsep=0.25em, topsep=0.30em}

% ---------- Tables ----------
\usepackage{booktabs}
\usepackage{array}
\usepackage{tabularx}
\renewcommand{\arraystretch}{1.15}

% ---------- Footnotes — keep tight, per-page ----------
\usepackage[bottom,hang,flushmargin]{footmisc}
\renewcommand{\footnotesize}{\fontsize{7.8pt}{9.2pt}\selectfont}
\setlength{\footnotesep}{0.7em}

% ---------- Blockquote styling ----------
\renewenvironment{quote}{\begin{list}{}{%
  \setlength{\leftmargin}{0.8em}%
  \setlength{\rightmargin}{0.3em}%
  \setlength{\topsep}{0.3em}%
  \setlength{\parsep}{0pt}%
}\item[]\itshape\small}{\end{list}}

% ---------- Hyperlinks (last) ----------
\usepackage[hidelinks,breaklinks=true]{hyperref}
\usepackage{xurl}                   % aggressive URL line-breaking for narrow pages
\hypersetup{
  pdfpagelayout=SinglePage,
  pdfdisplaydoctitle=true,
}

% ---------- Title block ----------
\makeatletter
\renewcommand{\maketitle}{%
  \begin{flushleft}%
    {\sffamily\Large\bfseries \@title\par}%
    \vspace{0.4em}%
    {\sffamily\small \@author\par}%
    \vspace{0.2em}%
    {\sffamily\footnotesize \@date\par}%
    \vspace{0.4em}%
    \hrule height 0.4pt%
    \vspace{0.8em}%
  \end{flushleft}%
}
\makeatother

% ---------- Convenience commands ----------
% Informal superscript reference marker for repeat citations.
\renewcommand{\cite}[1]{\textsuperscript{\textsf{\scriptsize #1}}}

%   \newcommand{\briefingtitle}{My Briefing}
%   \newcommand{\transcriptnum}{1}
%   \title{...}
%   \author{...}
%   \date{...}
%   \begin{document}\maketitle\thispagestyle{plain}
%   ... body ...
%   \end{document}
%
% ANNOTATION MARGIN:
%   The default right margin is 0.75in, which reserves space for
%   handwritten stylus annotations on e-ink readers. This gives a
%   text column of ~2.90in — a good balance between readability and
%   annotation space. For maximum words-per-page (no annotations),
%   change to `right=0.35in`. For a wider annotation column, use
%   `right=1.00in` (~2.65in text column, ~25% more pages).
% ==============================================================

\documentclass[10pt]{article}

% ---------- Geometry (Remarkable Pro Move, ~4in wide, 16:9) ----------
\usepackage[
  paperwidth=4.00in,
  paperheight=7.11in,
  top=0.40in,
  bottom=0.45in,
  left=0.35in,
  right=0.75in,         % annotation margin for stylus notes (0.35in for no margin)
  footskip=0.20in,
  headheight=12pt,
  headsep=0.10in
]{geometry}

% ---------- Fonts ----------
\usepackage[T1]{fontenc}
\usepackage[utf8]{inputenc}
\usepackage{tgpagella}              % Pagella for body (Latin-script serif)
\renewcommand{\familydefault}{\rmdefault}
\usepackage{tgadventor}             % Adventor (Latin-script sans) companion
\renewcommand{\sfdefault}{qag}      % set sans default to Adventor
\usepackage{microtype}              % cleaner line breaks and spacing

% ---------- Layout helpers ----------
\usepackage{setspace}
\setstretch{1.14}                   % mild leading for e-ink readability
\usepackage{parskip}                % blank line between paragraphs, no indent
\setlength{\parskip}{0.55em plus 0.15em minus 0.10em}

% narrow-column line breaking: allow TeX more slack before overfulling
\setlength{\emergencystretch}{3.5em}
\hyphenpenalty=50
\tolerance=2000

% ---------- Sectioning ----------
\usepackage{titlesec}
% Smaller headings — page is only 4in wide
\titleformat{\section}{\sffamily\large\bfseries}{\thesection}{0.6em}{}
\titleformat{\subsection}{\sffamily\normalsize\bfseries}{\thesubsection}{0.5em}{}
\titleformat{\subsubsection}{\sffamily\small\bfseries\itshape}{}{0em}{}
\titlespacing*{\section}{0pt}{1.4em}{0.6em}
\titlespacing*{\subsection}{0pt}{1.1em}{0.45em}
\titlespacing*{\subsubsection}{0pt}{0.8em}{0.30em}

% ---------- Page headers/footers ----------
\usepackage{fancyhdr}
\pagestyle{fancy}
\fancyhf{}
\renewcommand{\headrulewidth}{0.3pt}
\renewcommand{\footrulewidth}{0pt}
% The document MUST define \briefingtitle and \transcriptnum with
% \newcommand BEFORE \begin{document} (see example-body.tex).
\fancyhead[L]{\scriptsize\sffamily\MakeUppercase{\briefingtitle}}
\fancyhead[R]{\scriptsize\sffamily \transcriptnum}
\fancyfoot[C]{\small\thepage}

% Plain style for first page
\fancypagestyle{plain}{%
  \fancyhf{}%
  \renewcommand{\headrulewidth}{0pt}%
  \fancyfoot[C]{\small\thepage}%
}

% ---------- Lists ----------
\usepackage{enumitem}
\setlist{nosep, leftmargin=1.1em, itemsep=0.25em, topsep=0.30em}

% ---------- Tables ----------
\usepackage{booktabs}
\usepackage{array}
\usepackage{tabularx}
\renewcommand{\arraystretch}{1.15}

% ---------- Footnotes — keep tight, per-page ----------
\usepackage[bottom,hang,flushmargin]{footmisc}
\renewcommand{\footnotesize}{\fontsize{7.8pt}{9.2pt}\selectfont}
\setlength{\footnotesep}{0.7em}

% ---------- Blockquote styling ----------
\renewenvironment{quote}{\begin{list}{}{%
  \setlength{\leftmargin}{0.8em}%
  \setlength{\rightmargin}{0.3em}%
  \setlength{\topsep}{0.3em}%
  \setlength{\parsep}{0pt}%
}\item[]\itshape\small}{\end{list}}

% ---------- Hyperlinks (last) ----------
\usepackage[hidelinks,breaklinks=true]{hyperref}
\usepackage{xurl}                   % aggressive URL line-breaking for narrow pages
\hypersetup{
  pdfpagelayout=SinglePage,
  pdfdisplaydoctitle=true,
}

% ---------- Title block ----------
\makeatletter
\renewcommand{\maketitle}{%
  \begin{flushleft}%
    {\sffamily\Large\bfseries \@title\par}%
    \vspace{0.4em}%
    {\sffamily\small \@author\par}%
    \vspace{0.2em}%
    {\sffamily\footnotesize \@date\par}%
    \vspace{0.4em}%
    \hrule height 0.4pt%
    \vspace{0.8em}%
  \end{flushleft}%
}
\makeatother

% ---------- Convenience commands ----------
% Informal superscript reference marker for repeat citations.
\renewcommand{\cite}[1]{\textsuperscript{\textsf{\scriptsize #1}}}

%   \newcommand{\briefingtitle}{My Briefing}
%   \newcommand{\transcriptnum}{1}
%   \title{...}
%   \author{...}
%   \date{...}
%   \begin{document}\maketitle\thispagestyle{plain}
%   ... body ...
%   \end{document}
%
% ANNOTATION MARGIN:
%   The default right margin is 0.75in, which reserves space for
%   handwritten stylus annotations on e-ink readers. This gives a
%   text column of ~2.90in — a good balance between readability and
%   annotation space. For maximum words-per-page (no annotations),
%   change to `right=0.35in`. For a wider annotation column, use
%   `right=1.00in` (~2.65in text column, ~25% more pages).
% ==============================================================

\documentclass[10pt]{article}

% ---------- Geometry (Remarkable Pro Move, ~4in wide, 16:9) ----------
\usepackage[
  paperwidth=4.00in,
  paperheight=7.11in,
  top=0.40in,
  bottom=0.45in,
  left=0.35in,
  right=0.75in,         % annotation margin for stylus notes (0.35in for no margin)
  footskip=0.20in,
  headheight=12pt,
  headsep=0.10in
]{geometry}

% ---------- Fonts ----------
\usepackage[T1]{fontenc}
\usepackage[utf8]{inputenc}
\usepackage{tgpagella}              % Pagella for body (Latin-script serif)
\renewcommand{\familydefault}{\rmdefault}
\usepackage{tgadventor}             % Adventor (Latin-script sans) companion
\renewcommand{\sfdefault}{qag}      % set sans default to Adventor
\usepackage{microtype}              % cleaner line breaks and spacing

% ---------- Layout helpers ----------
\usepackage{setspace}
\setstretch{1.14}                   % mild leading for e-ink readability
\usepackage{parskip}                % blank line between paragraphs, no indent
\setlength{\parskip}{0.55em plus 0.15em minus 0.10em}

% narrow-column line breaking: allow TeX more slack before overfulling
\setlength{\emergencystretch}{3.5em}
\hyphenpenalty=50
\tolerance=2000

% ---------- Sectioning ----------
\usepackage{titlesec}
% Smaller headings — page is only 4in wide
\titleformat{\section}{\sffamily\large\bfseries}{\thesection}{0.6em}{}
\titleformat{\subsection}{\sffamily\normalsize\bfseries}{\thesubsection}{0.5em}{}
\titleformat{\subsubsection}{\sffamily\small\bfseries\itshape}{}{0em}{}
\titlespacing*{\section}{0pt}{1.4em}{0.6em}
\titlespacing*{\subsection}{0pt}{1.1em}{0.45em}
\titlespacing*{\subsubsection}{0pt}{0.8em}{0.30em}

% ---------- Page headers/footers ----------
\usepackage{fancyhdr}
\pagestyle{fancy}
\fancyhf{}
\renewcommand{\headrulewidth}{0.3pt}
\renewcommand{\footrulewidth}{0pt}
% The document MUST define \briefingtitle and \transcriptnum with
% \newcommand BEFORE \begin{document} (see example-body.tex).
\fancyhead[L]{\scriptsize\sffamily\MakeUppercase{\briefingtitle}}
\fancyhead[R]{\scriptsize\sffamily \transcriptnum}
\fancyfoot[C]{\small\thepage}

% Plain style for first page
\fancypagestyle{plain}{%
  \fancyhf{}%
  \renewcommand{\headrulewidth}{0pt}%
  \fancyfoot[C]{\small\thepage}%
}

% ---------- Lists ----------
\usepackage{enumitem}
\setlist{nosep, leftmargin=1.1em, itemsep=0.25em, topsep=0.30em}

% ---------- Tables ----------
\usepackage{booktabs}
\usepackage{array}
\usepackage{tabularx}
\renewcommand{\arraystretch}{1.15}

% ---------- Footnotes — keep tight, per-page ----------
\usepackage[bottom,hang,flushmargin]{footmisc}
\renewcommand{\footnotesize}{\fontsize{7.8pt}{9.2pt}\selectfont}
\setlength{\footnotesep}{0.7em}

% ---------- Blockquote styling ----------
\renewenvironment{quote}{\begin{list}{}{%
  \setlength{\leftmargin}{0.8em}%
  \setlength{\rightmargin}{0.3em}%
  \setlength{\topsep}{0.3em}%
  \setlength{\parsep}{0pt}%
}\item[]\itshape\small}{\end{list}}

% ---------- Hyperlinks (last) ----------
\usepackage[hidelinks,breaklinks=true]{hyperref}
\usepackage{xurl}                   % aggressive URL line-breaking for narrow pages
\hypersetup{
  pdfpagelayout=SinglePage,
  pdfdisplaydoctitle=true,
}

% ---------- Title block ----------
\makeatletter
\renewcommand{\maketitle}{%
  \begin{flushleft}%
    {\sffamily\Large\bfseries \@title\par}%
    \vspace{0.4em}%
    {\sffamily\small \@author\par}%
    \vspace{0.2em}%
    {\sffamily\footnotesize \@date\par}%
    \vspace{0.4em}%
    \hrule height 0.4pt%
    \vspace{0.8em}%
  \end{flushleft}%
}
\makeatother

% ---------- Convenience commands ----------
% Informal superscript reference marker for repeat citations.
\renewcommand{\cite}[1]{\textsuperscript{\textsf{\scriptsize #1}}}

%   \newcommand{\briefingtitle}{My Briefing}
%   \newcommand{\transcriptnum}{1}
%   \title{...}
%   \author{...}
%   \date{...}
%   \begin{document}\maketitle\thispagestyle{plain}
%   ... body ...
%   \end{document}
%
% ANNOTATION MARGIN:
%   The default right margin is 0.75in, which reserves space for
%   handwritten stylus annotations on e-ink readers. This gives a
%   text column of ~2.90in — a good balance between readability and
%   annotation space. For maximum words-per-page (no annotations),
%   change to `right=0.35in`. For a wider annotation column, use
%   `right=1.00in` (~2.65in text column, ~25% more pages).
% ==============================================================

\documentclass[10pt]{article}

% ---------- Geometry (Remarkable Pro Move, ~4in wide, 16:9) ----------
\usepackage[
  paperwidth=4.00in,
  paperheight=7.11in,
  top=0.40in,
  bottom=0.45in,
  left=0.35in,
  right=0.75in,         % annotation margin for stylus notes (0.35in for no margin)
  footskip=0.20in,
  headheight=12pt,
  headsep=0.10in
]{geometry}

% ---------- Fonts ----------
\usepackage[T1]{fontenc}
\usepackage[utf8]{inputenc}
\usepackage{tgpagella}              % Pagella for body (Latin-script serif)
\renewcommand{\familydefault}{\rmdefault}
\usepackage{tgadventor}             % Adventor (Latin-script sans) companion
\renewcommand{\sfdefault}{qag}      % set sans default to Adventor
\usepackage{microtype}              % cleaner line breaks and spacing

% ---------- Layout helpers ----------
\usepackage{setspace}
\setstretch{1.14}                   % mild leading for e-ink readability
\usepackage{parskip}                % blank line between paragraphs, no indent
\setlength{\parskip}{0.55em plus 0.15em minus 0.10em}

% narrow-column line breaking: allow TeX more slack before overfulling
\setlength{\emergencystretch}{3.5em}
\hyphenpenalty=50
\tolerance=2000

% ---------- Sectioning ----------
\usepackage{titlesec}
% Smaller headings — page is only 4in wide
\titleformat{\section}{\sffamily\large\bfseries}{\thesection}{0.6em}{}
\titleformat{\subsection}{\sffamily\normalsize\bfseries}{\thesubsection}{0.5em}{}
\titleformat{\subsubsection}{\sffamily\small\bfseries\itshape}{}{0em}{}
\titlespacing*{\section}{0pt}{1.4em}{0.6em}
\titlespacing*{\subsection}{0pt}{1.1em}{0.45em}
\titlespacing*{\subsubsection}{0pt}{0.8em}{0.30em}

% ---------- Page headers/footers ----------
\usepackage{fancyhdr}
\pagestyle{fancy}
\fancyhf{}
\renewcommand{\headrulewidth}{0.3pt}
\renewcommand{\footrulewidth}{0pt}
% The document MUST define \briefingtitle and \transcriptnum with
% \newcommand BEFORE \begin{document} (see example-body.tex).
\fancyhead[L]{\scriptsize\sffamily\MakeUppercase{\briefingtitle}}
\fancyhead[R]{\scriptsize\sffamily \transcriptnum}
\fancyfoot[C]{\small\thepage}

% Plain style for first page
\fancypagestyle{plain}{%
  \fancyhf{}%
  \renewcommand{\headrulewidth}{0pt}%
  \fancyfoot[C]{\small\thepage}%
}

% ---------- Lists ----------
\usepackage{enumitem}
\setlist{nosep, leftmargin=1.1em, itemsep=0.25em, topsep=0.30em}

% ---------- Tables ----------
\usepackage{booktabs}
\usepackage{array}
\usepackage{tabularx}
\renewcommand{\arraystretch}{1.15}

% ---------- Footnotes — keep tight, per-page ----------
\usepackage[bottom,hang,flushmargin]{footmisc}
\renewcommand{\footnotesize}{\fontsize{7.8pt}{9.2pt}\selectfont}
\setlength{\footnotesep}{0.7em}

% ---------- Blockquote styling ----------
\renewenvironment{quote}{\begin{list}{}{%
  \setlength{\leftmargin}{0.8em}%
  \setlength{\rightmargin}{0.3em}%
  \setlength{\topsep}{0.3em}%
  \setlength{\parsep}{0pt}%
}\item[]\itshape\small}{\end{list}}

% ---------- Hyperlinks (last) ----------
\usepackage[hidelinks,breaklinks=true]{hyperref}
\usepackage{xurl}                   % aggressive URL line-breaking for narrow pages
\hypersetup{
  pdfpagelayout=SinglePage,
  pdfdisplaydoctitle=true,
}

% ---------- Title block ----------
\makeatletter
\renewcommand{\maketitle}{%
  \begin{flushleft}%
    {\sffamily\Large\bfseries \@title\par}%
    \vspace{0.4em}%
    {\sffamily\small \@author\par}%
    \vspace{0.2em}%
    {\sffamily\footnotesize \@date\par}%
    \vspace{0.4em}%
    \hrule height 0.4pt%
    \vspace{0.8em}%
  \end{flushleft}%
}
\makeatother

% ---------- Convenience commands ----------
% Informal superscript reference marker for repeat citations.
\renewcommand{\cite}[1]{\textsuperscript{\textsf{\scriptsize #1}}}
